\documentclass[sigconf]{acmart}

\AtBeginDocument{%
  \providecommand\BibTeX{{%
    \normalfont B\kern-0.5em{\scshape i\kern-0.25em b}\kern-0.8em\TeX}}}

\begin{document}

\title{AutoFlow: Intelligent Workflow Automation via Large Language Models}

\author{Author Name}
\email{author@example.com}
\affiliation{%
  \institution{Institution}
  \city{City}
  \country{Country}
}

\begin{abstract}
Workflow automation traditionally requires rigid programming and significant manual effort. We present AutoFlow, an intelligent system that leverages Large Language Models to interpret natural language instructions and automatically generate, execute, and monitor complex workflows across multiple software interfaces.
\end{abstract}

\keywords{Workflow Automation, Large Language Models, Intelligent User Interfaces}

\maketitle

\section{Introduction}
Modern knowledge work involves navigating multiple applications and performing repetitive tasks. AutoFlow aims to bridge the gap between human intent and automated execution by providing a natural language interface for orchestrating complex, multi-step workflows.

\section{Methodology}
AutoFlow parses user instructions into a directed acyclic graph (DAG) of actionable nodes. Each node represents a distinct API call or UI interaction, dynamically generated by the underlying LLM.
\subsection{System Architecture}
The system consists of an Intent Parser, a Workflow Generator, and an Execution Engine with built-in error recovery mechanisms.

\section{Experiments}
\subsection{Evaluation Setup}
We evaluate AutoFlow on a suite of common office tasks, measuring task completion rate and time savings.
\subsection{Results}
% Placeholder results
AutoFlow successfully completes 85\% of evaluated workflows with zero human intervention, significantly reducing task completion time.

\section{Conclusion}
AutoFlow demonstrates the potential of LLMs to act as versatile intelligent agents for personal and enterprise workflow automation. 

\bibliographystyle{ACM-Reference-Format}
% \bibliography{bibliography}

\end{document}
